% 海洋行星（综述）
% license CCBYSA3
% type Wiki

本文根据 CC-BY-SA 协议转载翻译自维基百科\href{https://en.wikipedia.org/wiki/Ocean_world}{相关文章}。

\begin{figure}[ht]
\centering
\includegraphics[width=6cm]{./figures/76246827a4922672.png}
\caption{地球表面主要由海洋覆盖，海洋占据了地球表面的 $75\%$。因此，地球可以被视为一个水世界，但并非一个完全的海洋行星。} \label{fig_Ocwo_1}
\end{figure}
海洋世界（ocean world）、海洋行星（ocean planet）或水世界（water world）是一类行星或天然卫星，其水圈中包含大量以海洋形式存在的水，这些海洋可以位于表面之下，形成地下海洋，或直接存在于表面，甚至可能淹没所有陆地$^{[1][2][3][4]}$。“海洋世界”一词有时也被用于指代其海洋由其他流体或造海物质（thalassogen）构成的天体$^{[5]}$，例如熔岩（如木卫一 Io）、氨（与水形成共晶混合物，可能存在于土卫六 Titan 的内部海洋中），或碳氢化合物（如土卫六表面的情况，其可能是系外海洋中最常见的一类）$^{[6]}$。对地外海洋的研究被称为行星海洋学（planetary oceanography）。

地球是目前唯一已知在其表面存在液态水体的天文天体，尽管在木星的卫星欧罗巴（Europa）与盖尼米德（Ganymede），以及土星的卫星恩克拉多斯（Enceladus）与土卫六（Titan）上，人们推测其表面之下可能存在地下海洋$^{[7]}$。已经发现若干系外行星具备支持液态水存在的适宜条件$^{[8]}$。此外，在地球上也发现了大量地下水，主要以含水层的形式存在$^{[9]}$。对于系外行星而言，受限于当前技术，尚无法直接观测其表面的液态水，因此通常以大气中的水汽作为替代性指标$^{[10]}$。海洋世界的性质为理解其自身历史以及整个太阳系的形成与演化提供了重要线索；与此同时，它们在生命起源与维持方面的潜力也引起了广泛关注。

2020 年 6 月，NASA 科学家基于数学建模研究报告称，在银河系中拥有海洋的系外行星很可能是普遍存在的$^{[11][12]}$。
\subsection{概述}  
\subsubsection{定义}  
根据 Lunine 的定义，“海洋”被界定为“稳定存在、环绕整个天体的液态水体”$^{[13]}$。此外，“海洋世界”这一术语用于指代太阳系中那些拥有稳定、覆盖全球的液态水体的天体；与之相对，“海洋行星”和“水世界”则专指系外行星（即绕其他恒星运行的行星），其整体组成中含有较高质量分数的水$^{[13]}$。
\subsubsection{太阳系行星天体}
\begin{figure}[ht]
\centering
\includegraphics[width=8cm]{./figures/4b85d59dce5a977e.png}
\caption{土卫二（Enceladus）内部结构示意图} \label{fig_Ocwo_2}
\end{figure}
海洋世界因其可能孕育生命并维持生物活动的潜力而受到天体生物学家的高度关注$^{[4][3]}$。太阳系中被认为拥有地下海洋的主要卫星和矮行星尤为重要，因为与远在数光年之外、受当前技术所限难以直接探测的系外行星不同，这些天体可以通过空间探测器抵达并开展实地研究。除地球之外，太阳系中证据最为充分的水世界包括卡利斯托（Callisto）、恩克拉多斯（Enceladus）、欧罗巴（Europa）、盖尼米德（Ganymede）以及土卫六（Titan）$^{[3][14]}$。其中，欧罗巴和恩克拉多斯由于其外壳较薄且存在冰火山活动，被视为极具吸引力的探测目标。

太阳系中的其他天体也被视为可能拥有地下海洋的候选者，这一判断基于单一类型的观测证据或理论建模结果，包括天卫一（Ariel）$^{[14]}$、天卫三（Titania）$^{[15][16]}$、天卫二（Umbriel）$^{[17]}$、谷神星（Ceres）$^{[3]}$、土卫四（Dione）$^{[18]}$、土卫一（Mimas）$^{[19][20]}$、天卫五（Miranda）$^{[14]}$、天卫四（Oberon）$^{[4][21]}$、冥王星（Pluto）$^{[22]}$、海卫一（Triton）$^{[23]}$、厄里斯（Eris）$^{[4][24]}$以及鸟神星（Makemake）$^{[24]}$。
\subsubsection{系外行星}
\begin{figure}[ht]
\centering
\includegraphics[width=8cm]{./figures/d53b32c9e8bb7ae8.png}
\caption{一组含水、大小各异的系外行星，与地球进行对比（艺术家概念图；2018 年 8 月 17 日）$^{[25]}$} \label{fig_Ocwo_3}
\end{figure}
\begin{figure}[ht]
\centering
\includegraphics[width=8cm]{./figures/c7e845b1bfbd19f5.png}
\caption{以纯海洋世界作为过渡群体的系外行星族群示意图：它们位于气态巨行星与熔岩行星或岩石行星之间。} \label{fig_Ocwo_4}
\end{figure}
在太阳系之外，被描述为**候选海洋世界**的系外行星包括 GJ 1214 b $^{[26][27]}$、Kepler-22b、Kepler-62e、Kepler-62f $^{[28][29][30][31]}$，以及 Kepler-11 $^{[32]}$ 和 TRAPPIST-1 $^{[33][34]}$ 行星系统中的行星。

近年来，系外行星 TOI-1452 b、Kepler-138c 和 Kepler-138d 被发现其密度与其质量中有很大一部分由水构成的情形相一致 $^{[35][36]}$。此外，对高质量岩石行星 LHS 1140 b 的模型研究表明，其表面可能被一片深海所覆盖 $^{[37]}$。

尽管地球表面有 70.8\% 被水覆盖 $^{[38]}$，水在地球总质量中仅占约 0.05\%。在地外海洋世界中，海洋可能深且致密，以至于即便在高温条件下，巨大的压力也会使水转变为冰。在这类海洋的下部区域，数千巴的巨大压力可能促成奇异形态冰（如冰 V）地幔的形成 $^{[32]}$。这种冰并不一定像通常意义上的冰那样寒冷。若行星距离其恒星足够近，使水达到沸点，则水将进入超临界状态，不再具有清晰界定的表面 $^{[39]}$。即便是在较冷、以水为主的行星上，其大气层也可能比地球厚得多，并主要由水汽构成，从而产生极强的温室效应。此类行星必须足够小，无法保留厚重的氢和氦包层 $^{[40]}$，或者足够靠近其主星，以致这些轻元素被剥离 $^{[32]}$；否则，它们将演化成类似天王星和海王星那样的、更温暖版本的冰巨行星（此处尚需进一步证据）。
\subsection{历史}
20 世纪初的引力计算曾提出，木卫二（Europa）的组成可能富含水。1957 年，杰拉德·柯伊伯（Gerard Kuiper）通过地基观测证实了其表面主要由水冰构成。$^{[41]}$

在 20 世纪 70 年代行星探测任务之前，相关的重要理论基础工作已相继完成。1971 年，Lewis 指出，仅放射性衰变就可能足以在大型卫星中形成地下海洋，尤其是在存在氨（NH$_3$）的情况下。1979 年，Peale 和 Cassen 揭示了潮汐加热（亦称潮汐形变）在卫星演化与内部结构中的关键作用。$^{[3]}$
1992 年，人类首次确认探测到一颗系外行星。2003 年，Marc Kuchner 以及 2004 年 Alain Léger 等人提出，一小部分在雪线之外形成的冰质行星可能向内迁移至约 $1,\mathrm{AU}$ 的位置，其外层随后发生融化。$^{[42][43]}$

来自哈勃空间望远镜，以及先锋号、伽利略号、旅行者号、卡西尼–惠更斯号和新视野号等任务所积累的证据，强烈表明太阳系外侧的多个天体在隔热的冰壳之下蕴藏着内部液态水海洋。$^{[3][44]}$
与此同时，于 2009 年 3 月 7 日发射的开普勒空间望远镜已发现数千颗系外行星，其中约有 50 颗大小与地球相当，位于宜居带内或附近。$^{[45][46]}$

目前已探测到质量、尺寸和轨道各异的行星，这不仅展示了行星形成过程的多样性，也反映了行星在形成之后通过环恒星盘发生迁移的普遍现象。$^{[10]}$
截至 2025 年 10 月 30 日，已确认的系外行星共有 6,128 颗，分布在 4,584 个行星系统中，其中 1,017 个系统包含多于一颗行星。$^{[47]}$

2020 年 6 月，NASA 科学家基于数学建模研究报告称，银河系中拥有海洋的系外行星很可能相当普遍。$^{[11]}$

2022 年 8 月，距离地球约 99 光年的超级地球系外行星 TOI-1452 b 被凌日系外行星巡天卫星发现，该行星可能拥有深层海洋。$^{[35]}$
\subsection{形成}
\begin{figure}[ht]
\centering
\includegraphics[width=6cm]{./figures/00fa51d11bed7959.png}
\caption{阿塔卡马大型毫米/亚毫米阵列（ALMA）拍摄的 HL 金牛座（HL Tauri）原行星盘图像} \label{fig_Ocwo_5}
\end{figure}
在太阳系外侧形成的行星体，其初始状态通常类似彗星，按质量计大约由一半的水和一半的岩石组成，其整体密度低于岩质行星$^{[43]}$。在霜线附近形成的冰质行星和卫星应主要由 $\mathrm{H_2O}$ 和硅酸盐组成；而在更远区域形成的天体，则可以以水合物形式获得氨（$\mathrm{NH_3}$）和甲烷（$\mathrm{CH_4}$），并伴随一氧化碳（$\mathrm{CO}$）、氮气（$\mathrm{N_2}$）以及二氧化碳（$\mathrm{CO_2}$）$^{[48]}$。

在气态环星盘尚未消散之前形成的行星，会经历强烈的力矩作用，从而可能诱发其快速向内迁移进入宜居带，尤其对类地质量范围内的行星而言更为显著$^{[49][48]}$。由于水在岩浆中具有很高的溶解度，行星中相当大一部分水最初会被封存于地幔之中。随着行星逐渐冷却，地幔自下而上开始固结，大量水（约占地幔中水总量的 $60\%\sim99\%$）会发生析出，形成富含水汽的蒸汽大气层，并最终可能冷凝为海洋$^{[49]}$。海洋的形成需要行星内部发生分异过程，并且需要热i一个热源，这一热源可以来自放射性衰变、潮汐加热，或母体在早期阶段的高光度输出$^{[3]}$。然而，吸积完成之后的初始条件在理论上仍然缺乏充分约束。

在盘中外侧富水区域形成并随后向内迁移的行星，更有可能保有丰富的水储量$^{[50]}$。相反，在靠近其宿主恒星区域形成的行星，由于原始气体与尘埃盘的内侧通常被认为是高温且干燥的环境，因此其含水量往往较低。因而，如果在恒星附近发现一颗水世界，这将构成行星发生迁移并在原位之外形成的有力证据$^{[32]}$，因为在靠近恒星的区域缺乏足够的挥发性物质，难以支持原位形成$^{[2]}$。针对太阳系及系外行星系统形成过程的数值模拟表明，行星在形成过程中往往倾向于向内迁移（即朝向恒星）$^{[51][52][53]}$；在特定条件下，也可能发生向外迁移$^{[53]}$。向内迁移为冰质行星进入其冰层能够融化为液态水的轨道区域提供了可能性，从而使其转变为海洋行星。这一设想最早由 Marc Kuchner 于 2003 年在天文学文献中提出$^{[48]}$。
\subsection{结构}
冰质天体的内部结构通常通过对其整体密度、引力矩以及形状的测量来推断。确定天体的转动惯量有助于评估其是否经历过分异过程（即分离为岩石层与冰层）。在某些情况下，如果天体处于流体静力平衡状态（即在长时间尺度上表现为流体行为），则可以通过其形状或引力测量来反推转动惯量。尽管严格证明某一天体处于流体静力平衡状态极为困难，但通过综合利用形状与引力数据，仍可推导出其中的流体静力学贡献$^{[3]}$。用于探测内部海洋的具体技术包括磁感应测量、大地测量学、自转摆动、轴倾角、潮汐响应、雷达探测、成分证据以及表面形貌分析等方法$^{[3]}$。
\begin{figure}[ht]
\centering
\includegraphics[width=8cm]{./figures/1fed7f29913857c8.png}
\caption{木卫三（Ganymede）内部结构的艺术剖面示意图，其中液态水海洋被“夹”在两层冰层之间。各结构层按照真实比例绘制。} \label{fig_Ocwo_6}
\end{figure}
一种典型的冰质卫星由位于硅酸盐核心之上的水层构成。对于像土卫二（Enceladus）这样的小型卫星，海洋直接位于硅酸盐层之上、固态冰壳之下；而对于像木卫三（Ganymede）这样体量更大、富含冰的天体，其内部压力足够高，使得深处的冰转变为更高压相，从而有效形成一种“水夹层”结构，即液态海洋被夹在冰壳之间 $^{[3]}$。这两种情形之间的一个重要差异在于：在小型卫星中，海洋与硅酸盐直接接触，这可能为简单生命形式提供热液和化学能量及营养物质 $^{[3]}$。由于内部压力随深度变化，水世界模型中可能包含水的“蒸汽、液态、超流体、高压冰以及等离子体相”等多种状态 $^{[54]}$。其中部分固态水可能以冰 VII 的形式存在 $^{[55]}$。

维持地下海洋取决于内部加热速率与热量散失速率之间的对比，以及液体的凝固点 $^{[3]}$。因此，海洋的长期存在与潮汐加热机制密切相关。

较小的海洋行星通常具有更稀薄的大气层和更低的引力，因此液态水比在更大质量的海洋行星上更容易蒸发。数值模拟表明，质量小于一个地球质量的行星或卫星，仍可能在热液活动、放射性加热或潮汐形变的驱动下维持液态海洋 $^{[4]}$。当流体—岩石相互作用向深部脆性层传播较慢时，由蛇纹石化反应产生的热能可能成为小型海洋行星中热液活动的主要来源 $^{[4]}$。位于受潮汐形变影响的冰壳之下的全球性海洋，其动力学过程构成了一组重要而复杂的问题，目前才刚刚开始被系统研究。低温火山作用（cryovolcanism）发生的程度仍存在争议，因为水的密度比冰高约 8\%，在通常情况下不易喷发至表面 $^{[3]}$。尽管如此，旅行者 2 号、卡西尼–惠更斯号、伽利略号以及新视野号探测器的成像数据，已经在太阳系内多颗冰质天体表面揭示了低温火山活动的地貌特征。近期研究表明，在其内部海洋被表层冰覆盖的海洋行星上，低温火山作用可能与太阳系中的土卫二和木卫二（Europa）类似地发生 $^{[11][12]}$。

系外行星上的液态水海洋，其深度可能显著大于地球海洋的平均深度（约 3.7 km）$^{[56]}$。根据行星的引力和表面条件不同，系外行星的海洋深度可能达到地球的数十倍甚至数百倍。例如，对于一个表面温度约为 300 K 的行星，其液态水海洋深度可能在 30–500 km 之间，具体取决于行星的质量和组成 $^{[57]}$。
\subsection{大气模型}
\begin{figure}[ht]
\centering
\includegraphics[width=6cm]{./figures/bcebfd2db1b06d69.png}
\caption{具有氢大气层的 Hycean 行星的艺术想象图，这是一种大型海洋世界。} \label{fig_Ocwo_7}
\end{figure}
为了使地表水能够在长时间尺度上保持液态，一颗行星或卫星必须运行在宜居带（HZ）内，并拥有起保护作用的磁场 $^{[58][59][10]}$，同时具备足以维持足够大气压的引力 $^{[8]}$。如果行星的引力不足以支撑这一条件，那么其水最终将全部蒸发并逸散到外层空间。由电导流体层内部发电机机制维持的强行星磁层，有助于屏蔽恒星风造成的上层大气质量损失，并在长地质时间尺度上保留水分 $^{[58]}$。

行星大气可在行星形成过程中通过逸气作用产生，或从周围的原行星星云中被引力俘获。系外行星的表面温度主要由其大气中的温室气体（或其缺失）所控制，因此，大气可以通过向外辐射的红外辐射被探测到，这是因为温室气体会吸收并重新辐射来自宿主恒星的能量 $^{[10]}$。那些形成于富冰区域、随后向内迁移并进入过于靠近宿主恒星轨道的冰质行星，可能会形成厚重的蒸汽状大气，但即便其大气经历缓慢的流体动力学逃逸，仍可能在数十亿年尺度上保留其挥发性物质 $^{[42][48]}$。紫外光子不仅对生命具有生物学危害，还能够驱动快速的大气逃逸，从而导致行星大气的侵蚀 $^{[49][48]}$；水汽的光解以及氢和氧向太空的逃逸，可能使宜居带内的行星损失相当于多个地球海洋的水量，而这一过程无论是在能量受限还是扩散受限条件下均可发生 $^{[49]}$。水的损失量似乎与行星质量成正比，因为在扩散受限条件下，氢的逃逸通量与行星表面引力成正相关。

在失控温室效应期间，水汽会到达平流层，在那里容易被紫外辐射光解。紫外辐射对高层大气的加热随后可能驱动一种流体动力学风，将氢（以及可能的一部分氧）带入太空，从而导致行星地表水的不可逆损失、地表的氧化，以及大气中氧的潜在积累 $^{[49]}$。某一特定行星大气的最终命运，在很大程度上取决于极紫外辐射通量、失控温室阶段的持续时间、初始水含量，以及氧被地表吸收的速率 $^{[49]}$。富含挥发物的行星在年轻恒星以及 M 型恒星的宜居带中可能更为常见 $^{[48]}$。

科学家还提出了“Hycean 行星”的概念，即具有主要由氢构成的厚重大气层的海洋行星。这类行星在其恒星周围可能拥有较宽的轨道范围，在其中可以维持液态水。然而，这些模型往往基于相对简化的行星大气假设。更为复杂的研究表明，氢与恒星光在不同波长下的相互作用方式，与氮和氧等较重元素存在显著差异。如果这样一颗行星，其大气压为地球的 10–20 倍，并位于距离恒星 1 个天文单位（AU）处，那么其水体将会沸腾。最新研究因此将此类世界的宜居带位置修正为约 3.85 AU；若其大气压与地球相近，则宜居带位置约为 1.6 AU $^{[60]}$。
\subsubsection{成分模型}
在研究系外行星的表面及其大气时面临诸多挑战，因为云层覆盖会影响大气温度、大气结构以及光谱特征的可观测性 $^{[61]}$。然而，位于宜居带（HZ）内、由大量水构成的行星，预计其表面与大气将表现出独特的地球物理与地球化学特征 $^{[61]}$。例如，对于系外行星 Kepler-62e 与 Kepler-62f，其外表面可能是液态海洋、蒸汽大气，或完全被 Ice I（常压水冰）覆盖，这取决于它们在宜居带内的具体轨道位置以及温室效应的强度。还有多种表层与内部过程会影响大气成分，其中包括但不限于：用于溶解 CO$_2$ 的海洋比例、大气相对湿度、行星表面与内部的氧化还原状态、海洋酸度、行星反照率以及表面引力 $^{[10][62]}$。

大气结构以及由此确定的宜居带边界，取决于行星大气的密度：对于低质量行星，宜居带向外移动，而对于高质量行星，则向内移动 $^{[61]}$。理论研究与数值模型均表明，宜居带内水世界的行星大气成分，与同时具有陆地和海洋的行星相比，不应存在显著差异 $^{[61]}$。在建模时，通常假设组成水世界的冰质微行星体的初始成分与彗星相似，即主要由水（H$_2$O）构成，并含有一定比例的氨（NH$_3$）和二氧化碳（CO$_2$）$^{[61]}$。这种类似彗星的初始冰成分，会导致其大气模型中约 90\% 为 H$_2$O、5\% 为 NH$_3$、5\% 为 CO$_2$ $^{[61][63]}$。

针对 Kepler-62f 的大气模型显示，需要约 1.6–5 bar 的 CO$_2$ 大气压，才能将其表面温度加热至冰点以上，对应的表面压力约为地球的 0.56–1.32 倍 $^{[61]}$。
\subsection{海洋学}
研究表明，在土卫二（Enceladus）、土卫六（Titan）、木卫三（Ganymede）以及木卫二（Europa）中，可能存在强烈的海洋环流 $^{[64][65]}$。在土卫二中，由冰壳厚度推断出的海洋热通量表明，暖水在极区上涌，而较冷的水体在低纬度区域下沉 $^{[66][67]}$。模型预测木卫二存在赤道区域的暖水上涌，并在低纬度地区具有更强的热量传输 $^{[64]}$。在全球尺度上，其洋流被组织成三个纬向环流单元和两个赤道环流单元，将内部热量对流输送至表面，尤其在赤道区域最为显著 $^{[68][69][70]}$。相比之下，土卫六和木卫三被假设表现为近似非旋转体系，因此不具有一致、有序的热量传输模式 $^{[64]}$。
\subsection{天体生物学}
海洋世界或海洋行星的特征为其自身历史，以及整个太阳系的形成与演化提供了重要线索。除此之外，它们是否能够孕育并维持生命也备受关注。就我们所知，生命的存在需要液态水、能量来源以及营养物质，而在这些天体中的一部分，这三项关键条件都有可能得到满足 $^{[3]}$，从而在地质时间尺度上维持简单的生物活动 $^{[3][4]}$。2018 年 8 月，研究人员报告称，水世界可能具备支持生命存在的能力 $^{[71][72]}$。

然而，若一颗海洋世界的表面被完全覆盖的液态水所占据，其对类地生命的宜居性将受到限制；若全球海洋与下方岩石地幔之间还隔着一层高压固态冰层，则这种限制将更加严重 $^{[73][74]}$。对一种假想的、表面覆盖相当于五个地球海洋水量的海洋世界的数值模拟表明，其海水中将缺乏足够的磷及其他营养元素，难以支持类似地球上浮游生物那样、能够产生氧气的海洋生物演化。在地球上，磷主要通过雨水冲刷裸露陆地岩石进入海洋，而这一机制在完全被海洋覆盖的行星上无法发挥作用。进一步的模拟显示，若海洋行星拥有相当于 50 个地球海洋的水量，其海底压力将极其巨大，以至于行星内部难以维持板块构造活动，从而无法通过火山作用提供适合类地生命的化学环境 $^{[75]}$。

另一方面，像木卫二（Europa）和土卫二（Enceladus）这样的小型天体，被认为是尤为有利的宜居环境，因为理论上它们的海洋几乎可以肯定会与下方的硅酸盐核心直接接触，而后者可能同时提供热量以及对生命至关重要的化学元素 $^{[3]}$。这些天体的表面地质活动还可能将生物学上重要的“构建模块”输送至海洋中，例如来自彗星的有机分子，或由太阳紫外辐射作用于甲烷、乙烷等简单有机化合物并常与氮结合而形成的托林（tholins）$^{[76]}$。
\subsubsection{氧}
分子氧（$O_2$）既可以通过地球物理过程产生，也可以作为生命体进行光合作用的副产物出现。因此，尽管 $O_2$ 的存在令人鼓舞，但它并不是一种可靠的生物标志 $^{[39][49][77][10]}$。事实上，大气中 $O_2$ 浓度很高的行星甚至可能并不宜居 $^{[49]}$。在富含氧的大气环境中，生命的起源可能更加困难，因为早期生物依赖于多种含氢化合物参与的氧化还原反应所释放的自由能；而在 $O_2$ 丰富的行星上，生物体必须与氧竞争这种自由能 $^{[49]}$。
\subsection{另见}
\begin{itemize}
\item 环恒星宜居带——行星可能拥有液态地表水的轨道区域
\item 沙漠行星——表面以沙漠为主的岩石行星
\item 类地行星——环境与地球相似的行星
\item 地外液态水——在地球之外自然存在的液态水
\item 冰行星——表面被冰覆盖的行星
\item 系外液态水候选天体列表——地球之外可能存在液态水的天体
\item 海洋 § 地外海洋——覆盖地球大部分表面的咸水水体
\item 泛大洋（Panthalassa）——围绕盘古大陆的史前超级海洋
\item 《Subnautica》——2018 年电子游戏
\item TOI-1452 b——环绕 TOI-1452 的一颗超级地球
\end{itemize}
面向外太阳系水世界的天体生物学任务构想：

\begin{itemize}
\item Enceladus Explorer（土卫二探测器）——规划中的行星际轨道器与着陆器任务
\item Enceladus Life Finder（ELF，土卫二生命探测器）——美国国家航空航天局（NASA）提出的土星卫星探测任务
\item Europa Lander（木卫二着陆器）——已取消的 NASA 木卫二着陆任务
\item xplorer of Enceladus and Titan（E2T，土卫二与土卫六探测器）——NASA/ESA 提出的土星卫星探测概念
\item Journey to Enceladus and Titan（JET，前往土卫二与土卫六）——提出的空间探测任务
\item Jupiter Icy Moons Explorer（JUICE，木星冰卫星探测器）——自 2023 年起执行的欧洲空间局（ESA）木星及其卫星探测任务
\item Laplace-P——拟议中的俄罗斯航天器任务，用于研究木星卫星系统并在木卫三着陆
\item Life Investigation For Enceladus（LIFE，土卫二生命调查）——提出的天体生物学任务
\item Oceanus——2017 年提出的 NASA 海王星卫星 Triton（海卫一）轨道器探测任务
\item Testing the Habitability of Enceladus's Ocean（THEO，土卫二海洋宜居性测试）——土卫二轨道器任务
\item Titan Lake In-situ Sampling Propelled Explorer（TALISE，土卫六湖泊原位取样推进式探测器）——提出的空间探测任务
\item Titan Mare Explorer（TiME，土卫六海洋探测器）——提出的着陆器航天器设计方案
\item Triton Hopper（海卫一跳跃式探测器）——提出的 NASA 海卫一着陆探测任务
\end{itemize}
\subsection{参考文献}
\begin{enumerate}
\item “Ocean planet definition/meaning”（海洋行星的定义/含义）。Omnilexica，2017 年 10 月 1 日；2017 年 10 月 2 日存档。2017 年 10 月 1 日检索。海洋行星是一种假想的行星类型，其相当大一部分质量由水构成。这类行星的表面将被深达数百千米的水海完全覆盖，远深于地球海洋。
\item Adams, E. R.; Seager, S.; Elkins-Tanton, L.（2008 年 2 月 1 日）。“Ocean Planet or Thick Atmosphere: On the Mass–Radius Relationship for Solid Exoplanets with Massive Atmospheres”（海洋行星还是厚重大气层：具有巨大大气的固态系外行星的质量—半径关系）。《天体物理学杂志》，673(2)：1160–1164。arXiv:0710.4941。Bibcode:2008ApJ...673.1160A。doi:10.1086/524925。具有给定质量和半径的行星，可能含有大量水冰（即所谓的海洋行星），也可能拥有大型岩铁核心并伴随一定量的氢和/或氦。
\item Nimmo, F.; Pappalardo, R. T.（2016 年 8 月 8 日）。“Ocean worlds in the outer solar system”（外太阳系中的海洋世界）。《地球物理研究杂志》，121(8)：1378。Bibcode:2016JGRE..121.1378N。doi:10.1002/2016JE005081。
\item Vance, Steve; Harnmeijer, Jelte; Kimura, Jun; Hussmann, Hauke; Brown, J. Michael（2007）。“Hydrothermal Systems in Small Ocean Planets”（小型海洋行星中的热液系统）。《天体生物学》，7(6)：987–1005。Bibcode:2007AsBio...7..987V。doi:10.1089/ast.2007.0075。PMID 18163874。
\item 《Ocean Worlds: The story of seas on Earth and other planets》（《海洋世界：地球及其他行星上海洋的故事》），Jan Zalasiewicz 与 Mark Williams 著，牛津大学出版社，2014 年 10 月 23 日。ISBN 019165356X，9780191653568。
\item Ballesteros, F. J.; Fernandez-Soto, A.; Martinez, V. J.（2019）。“Diving into Exoplanets: Are Water Seas the Most Common?”（深入系外行星：水海是否最为常见？）。《天体生物学》，19(5)：642–654。doi:10.1089/ast.2017.1720。hdl:10261/213115。PMID 30789285。
\item “Ocean Worlds: Water in the Solar System and Beyond – NASA Science”（《海洋世界：太阳系内外的水》），NASA 科学，2023 年 7 月 22 日。
\item “Are there oceans on other planets?”（其他行星上是否存在海洋？）。美国国家海洋和大气管理局（NOAA），2017 年 7 月 6 日。2017-10-03 检索。
\item “Aquifers and Groundwater | U.S. Geological Survey”（含水层与地下水 | 美国地质调查局）。USGS，2018 年 10 月 9 日。2023-05-02 检索。
\item Seager, Sara（2013）。“Exoplanet Habitability”（系外行星宜居性）。《科学》，340：577–581。Bibcode:2013Sci...340..577S。doi:10.1126/science.1232226。PMID 23641111。
\item Shekhtman, Lonnie 等（2020 年 6 月 18 日）。“Are Planets with Oceans Common in the Galaxy? It’s Likely, NASA Scientists Find”（银河系中拥有海洋的行星是否常见？NASA 科学家认为很可能）。NASA。2020 年 6 月 20 日检索。
\item Quick, Lynnae C.; Roberge, Aki; Barr Mlinar, Amy; Hedman, Matthew M.（2020-06-18）。“Forecasting Rates of Volcanic Activity on Terrestrial Exoplanets and Implications for Cryovolcanic Activity on Extrasolar Ocean Worlds”（类地系外行星火山活动速率预测及其对系外海洋世界低温火山活动的启示）。《太平洋天文学会出版物》，132(1014)：084402。Bibcode:2020PASP..132h4402Q。doi:10.1088/1538-3873/ab9504。
\item Lunine, Jonathan I.（2017 年 2 月）。“Ocean worlds exploration”（海洋世界探测）。《航天学报》，131：123–130。Bibcode:2017AcAau.131..123L。doi:10.1016/j.actaastro.2016.11.017。
\item Hendrix, Amanda R. 等（2019）。“The NASA Roadmap to Ocean Worlds”（NASA 海洋世界路线图）。《天体生物学》，19(1)：1–27。Bibcode:2019AsBio..19....1H。doi:10.1089/ast.2018.1955。PMC 6338575。PMID 30346215。
\item “The Weather Channel – Ashburn, VA, United States”（天气频道——美国弗吉尼亚州阿什本）。
\item “New Study of Uranus’ Large Moons Shows 4 May Hold Water – NASA”（新研究显示天王星四颗大型卫星可能含水），NASA，2023 年 5 月 4 日。
\item “Uranus’ 4 biggest moons may have buried oceans of salty water”（天王星四大卫星可能拥有埋藏的咸水海洋）。Space.com，2023 年 5 月 5 日。
\item Zannoni, Marco; Hemingway, Douglas; Gomez Casajus, Luis; Tortora, Paolo（2020 年 7 月）。“The gravity field and interior structure of Dione”（土卫四的引力场与内部结构）。《伊卡洛斯》，345：113713。arXiv:1908.07284。Bibcode:2020Icar..34513713Z。doi:10.1016/j.icarus.2020.113713。
\item Ocean Worlds（海洋世界）。NASA 喷气推进实验室（JPL）。
\item Ocean Worlds Exploration Program（海洋世界探测计划）。NASA。
\item Hussmann, Hauke; Sohl, Frank; Spohn, Tilman（2006 年 11 月）。“Subsurface oceans and deep interiors of medium-sized outer planet satellites and large trans-neptunian objects”（中等大小外行星卫星与大型海王星外天体的地下海洋与深部内部）。《伊卡洛斯》，185(1)：258–273。Bibcode:2006Icar..185..258H。doi:10.1016/j.icarus.2006.06.005。
\item Johnson, Brandon C. 等（2016 年 10 月）。“Formation of the Sputnik Planum basin and the thickness of Pluto’s subsurface ocean”（冥王星斯普特尼克平原盆地的形成及其地下海洋厚度）。《地球物理研究快报》，43(19)：10,068–10,077。Bibcode:2016GeoRL..4310068J。doi:10.1002/2016GL070694。
\item Schenk, Paul 等（2021 年 9 月）。“Triton: Topography and Geology of a Probable Ocean World with Comparison to Pluto and Charon”（海卫一：一个可能的海洋世界的地形与地质，并与冥王星和卡戎比较）。《遥感》，13(17)：3476。Bibcode:2021RemS...13.3476S。doi:10.3390/rs13173476。hdl:10150/661940。
\item Glein, Christopher R. 等（2024 年 4 月）。“Moderate D/H ratios in methane ice on Eris and Makemake as evidence of hydrothermal or metamorphic processes in their interiors: Geochemical analysis”（阋神星与鸟神星甲烷冰中适中的 D/H 比作为其内部热液或变质过程的证据：地球化学分析）。《伊卡洛斯》，412：115999。arXiv:2309.05549。Bibcode:2024Icar..41215999G。doi:10.1016/j.icarus.2024.115999。
\item “Water-worlds are common: Exoplanets may contain vast amounts of water”（水世界很常见：系外行星可能含有大量水）。Phys.org，2018 年 8 月 17 日。
\item Charbonneau, David 等（2009）。“A super-Earth transiting a nearby low-mass star”（一颗凌日的近邻低质量恒星周围的超级地球）。《自然》，462：891–894。arXiv:0912.3229。Bibcode:2009Natur.462..891C。doi:10.1038/nature08679。PMID 20016595。
\item Kuchner, Seager; Hier-Majumder, M.; Militzer, C. A.（2007）。“Mass–radius relationships for solid exoplanets”（固态系外行星的质量—半径关系）。《天体物理学杂志》，669(2)：1279–1297。arXiv:0707.2895。Bibcode:2007ApJ...669.1279S。doi:10.1086/521346。
\item Water Worlds and Ocean Planets（《水世界与海洋行星》）。Sol Company，2012 年。
\item Rincon, Paul（2011 年 12 月 5 日）。“A home from home: Five planets that could host life”（似曾相识的家园：五颗可能孕育生命的行星）。BBC 新闻。2016 年 11 月 26 日检索。
\item D’Angelo, G.; Bodenheimer, P.（2016）。“In Situ and Ex Situ Formation Models of Kepler 11 Planets”（Kepler-11 行星的原位与异地形成模型）。《天体物理学杂志》，828(1)：33。arXiv:1606.08088。Bibcode:2016ApJ...828...33D。doi:10.3847/0004-637X/828/1/33。
\item Bourrier, Vincent; de Wit, Julien; Jäger, Mathias（2017 年 8 月 31 日）。“Hubble delivers first hints of possible water content of TRAPPIST-1 planets”（哈勃望远镜首次提供 TRAPPIST-1 行星可能含水的线索）。SpaceTelescope.org。2017 年 9 月 4 日检索。
\item PTI（2017 年 9 月 4 日）。“First evidence of water found on TRAPPIST-1 planets”（在 TRAPPIST-1 行星上发现水的首个证据）。《印度快报》。
\item Cadieux, Charles; Doyon, René 等（2022 年 9 月）。“TOI-1452 b: SPIRou and TESS Reveal a Super-Earth in a Temperate Orbit Transiting an M4 Dwarf”（TOI-1452 b：SPIRou 与 TESS 揭示的一颗绕 M4 型矮星运行、处于温和轨道的超级地球）。《天文学杂志》，164(3)：96。arXiv:2208.06333。Bibcode:2022AJ....164...96C。doi:10.3847/1538-3881/ac7cea。
\item Piaulet, Caroline; Benneke, Björn 等（2022 年 12 月 15 日）。“Evidence for the volatile-rich composition of a 1.5-Earth-radius planet”（一颗半径为 1.5 个地球半径行星富含挥发物成分的证据）。《自然·天文学》，7：206–222。arXiv:2212.08477。Bibcode:2023NatAs...7..206P。doi:10.1038/s41550-022-01835-4。
\item Lillo-Box, J.; Figueira, P. 等（2020 年 10 月）。“Planetary system LHS 1140 revisited with ESPRESSO and TESS”（利用 ESPRESSO 与 TESS 重新研究行星系统 LHS 1140）。《天文学与天体物理学》，642：A121。arXiv:2010.06928。Bibcode:2020A&A...642A.121L。doi:10.1051/0004-6361/202038922。hdl:10261/234118。
\item Pidwirny, M.（2006）。“Surface area of our planet covered by oceans and continents”（地球表面被海洋与大陆覆盖的面积）。英属哥伦比亚大学奥肯那根分校。2016 年 5 月 13 日检索。
\item Léger, Alain（2004）。“A New Family of Planets? ‘Ocean Planets’”（一种新的行星家族？“海洋行星”）。《伊卡洛斯》，169(2)：499–504。arXiv:astro-ph/0308324。Bibcode:2004Icar..169..499L。doi:10.1016/j.icarus.2004.01.001。
\end{enumerate}